\Article{247}
\textsc{Theorem.}\quad\emph{The product of two similar periods can be formed, independently of the number $p$, by adding similar periods and given numbers.}

(We call periods similar when they contain equally many terms, or when the number $e$ is the same.)

Example. Let $n=7$. The product of the periods $(1)+(6)$ and $(2)+(5)$ is, because $(a)\times(b)=(a+b)$, $(3)+(6)+(8)+(11)$, or it consists of the periods $(3)+(4)$ and $(1)+(6)$.

\emph{Proof.}\quad Let $\frac{\lambda}{e}=f$, and let the given periods be $(e\ast\alpha)$ and $(e\ast\beta)$, or the sums
\[
\begin{aligned}
&(\rho^\alpha)+(\rho^{\alpha+e})+(\rho^{\alpha+2e})+\ldots+(\rho^{\alpha+(f-1)e})\ \dotsfive\ P\\
&(\rho^\beta)+(\rho^{\beta+e})+(\rho^{\beta+2e})+\ldots+(\rho^{\beta+(f-1)e})\ \dotsfive\ Q
\end{aligned}
\]

The product $PQ$ consists of $f^2$ terms. Arrange them as follows. Form $f$ rows, each consisting of $f$ terms. Let the first contain the product of $P$ with $(\rho^\beta)$, the second the product $P\pdot(\rho^{\beta+e})$, and so on, and so on. In the first row let the product arising from the term $(\rho^\alpha)$ occupy the first place, that from $(\rho^{\alpha+e})$ the second, and the others follow in succession. In the second row assign the first place to the product arising from $(\rho^{\alpha+e})$, the second to that from $(\rho^{\alpha+2e})$, and so on, and finally the last place to that from $(\rho^\alpha)$. Let the third begin with the product arising from $(\rho^{\alpha+2e})$, and continue in this way, following the product from the last term with the products from the first and second terms, and so on, and so on. Thus, denoting the successive terms of the period $P$ by $1,2,3\ldots z$ and those of the period $Q$ by $\mathrm{I},\mathrm{II},\mathrm{III},\dotsfour Z$, arrange the terms of $PQ$ as
\[
\begin{aligned}
&1\pdot\mathrm{I}+2\pdot\mathrm{I}+3\pdot\mathrm{I}+4\pdot\mathrm{I}+\dotstwo+z\pdot\mathrm{I}\\
&2\pdot\mathrm{II}+3\pdot\mathrm{II}+4\pdot\mathrm{II}+\dotsfour+1\pdot\mathrm{II}\\
&3\pdot\mathrm{III}+4\pdot\mathrm{III}+\ldots+1\pdot\mathrm{III}+2\pdot\mathrm{III}\\
&\qquad\text{etc.}\quad\text{etc.}
\end{aligned}
\]

Now collect all terms occupying the same place in the individual rows into $f$ groups; and I say
\SourceLeaf{216}
\par\noindent $1^\circ$ if any term is $\equiv1$, all the other terms in the same group will also be $\equiv1$

\noindent $2^\circ$ every group in which no term is $\equiv1$ forms a period. --- Once these statements have been proved, we will evidently have achieved our purpose.

\Needspace{12\baselineskip}The general form of such a group is
\[
\begin{gathered}
(\rho^{\alpha+ke}+\rho^\beta),\quad(\rho^{\alpha+(k+1)e}+\rho^{\beta+e}),\quad(\rho^{\alpha+(k+2)e}+\rho^{\beta+2e}),\\
\ldots\ (\rho^{\alpha+(k+f-1)e}+\rho^{\beta+(f-1)e})
\end{gathered}
\]
Indeed, in place of $\rho^{\alpha+(k-1)e}$ one can also write $\rho^{\alpha+(k+f-1)e}$, since $ef=\lambda$ and $\rho^\lambda\equiv1\;(\mathrm{mod.}\,n)$, and similarly for the preceding terms. Put $\rho^{\alpha+ke}+\rho^\beta\equiv\rho^\chi\;(\mathrm{mod.}\,n)$, which is permissible unless $\rho^{\alpha+ke}+\rho^\beta$ happens to be divisible by $n$\begingroup\renewcommand{\thefootnote}{*)}\footnote{The proposition will have to be expressed somewhat differently if $n$ denotes a power of a prime in general; when it is a prime number, however, nothing needs changing.}\endgroup. The group can then be represented as $(\rho^\chi)$, $(\rho^{\chi+e})$, $(\rho^{\chi+2e})\ldots(\rho^{\chi+(f-1)e})$, which is evidently the period $(e\ast\chi)$. But if $\rho^{\alpha+ke}+\rho^\beta$ is divisible by $n$, all terms of the group will be $\equiv(0)$, that is, $\equiv1$. Q.\,E.\,D.

\emph{Note.} This proof also gives a very easy method of expanding the product. We shall give another below which lacks this advantage, but seems worthy of consideration for its simplicity.

\Article{248}
We call all the smaller periods that make up a larger period a system of periods. Thus this name will be given to the periods
\[
(ef\ast\alpha),\quad(ef\ast f+\alpha),\quad(ef\ast2f+\alpha)\ \dotsfour\ (ef\ast(e-1)f+\alpha)
\]
that make up the period $(f\ast\alpha)$. It will be \emph{properly ordered} if the numbers placed after the sign $\ast$, here $\alpha$, $f+\alpha$, $2f+\alpha$, follow an arithmetic progression with difference $f$. Finally, systems are \emph{similar} if both their smaller and their larger periods are similar.

\textsc{Theorem.} \emph{If the periods of two similar, properly ordered systems are multiplied together, namely the first by the first, the second by the second, the third by the third, and so on, the sum of all the products can be formed from given numbers and periods similar to the larger period.}

\emph{Proof.} Let the systems be
\[
\begin{gathered}
(ef\ast\alpha),\quad(ef\ast\alpha+f),\quad(ef\ast\alpha+2f)\ldots\\
(ef\ast\beta),\quad(ef\ast\beta+f),\quad(ef\ast\beta+2f)\ldots
\end{gathered}
\]

\SourceLeaf{217}
The products of the individual periods of the first system with the corresponding periods of the second consist, by the preceding section, of integers and similar periods. A little attention to how these periods arise shows that, if $(ef\ast\alpha)\times(ef\ast\beta)$ consists of an integer $N$ and periods $(ef\ast A)$, $(ef\ast B)$, $(ef\ast C)$, and so on, the products are then formed as follows:
\[
\begin{aligned}
&(ef\ast\alpha+f)\times(ef\ast\beta+f)\text{ from }N\text{ and periods}\\
&\hspace{2em}(ef\ast A+f),\quad(ef\ast B+f),\quad(ef\ast C+f)\text{ etc.}\\
&(ef\ast\alpha+2f)\times(ef\ast\beta+2f)\text{ from }N\text{ and periods}\\
&\hspace{2em}(ef\ast A+2f),\quad(ef\ast B+2f),\quad(ef\ast C+2f)\text{ etc.}
\end{aligned}
\]
and in general
\[
\begin{aligned}
&(ef\ast\alpha+\mu f)\times(ef\ast\beta+\mu f)\text{ from }N\text{ and periods}\\
&\hspace{2em}(ef\ast A+\mu f),\quad(ef\ast B+\mu f),\quad(ef\ast C+\mu f)\text{ etc.}
\end{aligned}
\]
It follows directly that the sum of all the periods will be
\[
eN+(f\ast A)+(f\ast B)+(f\ast C)\text{ etc.}\qquad\text{Q.\,E.\,D.}
\]
This proof also supplies a method for finding that sum.

\Article{249}
It is easy to make this theorem still more general. Namely, if there are any number of similar, properly ordered systems, and the products of all the first periods, all the second periods, and so on, are formed, the sum of all these products consists of numbers and larger periods. If all these systems are taken to be equal, the sum of powers of all the periods, whatever the exponent, will consist of numbers and periods similar to the larger one. The aim of these results is now apparent. Let $\lambda=efgh\ldots$. Divide all the primitive roots into $e$ periods $A,A',A''$, and so on, each of these again into $f$: $B,B',B''$, and so on, and each of these into $g$: $C,C',C''$, and so on. The sum of all the periods is given: it is $\equiv-1$. But according to what we have just said, we shall also have
\[
\begin{gathered}
(A)^2+(A')^2+(A'')^2+(A''')^2+\text{etc.}\\
(A)^3+(A')^3+(A'')^3+(A''')^3+\text{etc.}\\
\text{etc. etc.}
\end{gathered}
\]

Hence, by section 244, a congruence of degree $e$ can be found whose roots are $A,A',A''$, and so on. Taking these as known, let each period be divided into smaller ones:
\SourceLeaf{218}
\[
\begin{array}{rclll}
A&\text{into}&B,&B',&B''\ldots\\
A'&\text{into}&B^{(n)},&B^{(n+1)},&B^{(n+2)}\ldots\\
&\text{etc.}
\end{array}
\]
Thus $B+B'+B''+\ldots\equiv A$ is given. But
\[
\begin{gathered}
(B)^2+(B')^2+(B'')^2+\ldots\\
(B)^3+(B')^3+(B'')^3+\ldots\\
\text{etc.}
\end{gathered}
\]
consists of units and the periods $A,A',A''$, and so on. Therefore $B,B',B''$, and so on, will be given by a congruence of degree $f$, from which they can be found; and similarly the periods making up $A',A''$, and so on, can be determined. Anyone will now see that every period can be subdivided into smaller ones by precisely the same method, until the roots themselves are reached.

\Article{250}
But a difficulty occurs in applying these rules, and we must remove it. Since each congruence has several roots, we must decide which label to assign to each, so that they can be properly distinguished from one another. The labels of the periods depend on the numbers $r,\rho$, which can be chosen at will, so their labelling must also involve something arbitrary. The number $\rho$ must, however, be fixed at the outset. The main feature of our method is that we derive smaller periods from larger ones. But this cannot be done without the proper order of the periods, which we have obtained through their \emph{labels}. We must therefore ensure that all periods are distinguished by their labels as soon as they have been found.

Let the period $A$ be denoted by $(e\ast\alpha)$ and divided into $f$ periods $B,C,D$, and so on, which must be labelled. Each is clearly contained in a form $(ef\ast ke+\alpha)$ of this kind. But I say that, for one of them, $B$, the number $k$ can be chosen at will, and the placement of the others derived from it.

Let $R$ be any primitive root of the congruence $x^n\equiv1$, and let $B$ consist of the terms $R^\mu+R^\nu+\text{etc.}$. Put $\frac1\mu\rho^{ke+\alpha}\equiv\frac1\upsilon\;(\mathrm{mod.}\,n)$, and, since the value of $r$ is arbitrary, provided only that $A$ receives the label $(e\ast\alpha)$, which evidently happens of itself, put $r\equiv R^\upsilon\;(\mathrm{mod.}\,p)$. Then the first term of $B$ will be $r^{\rho^{ke+\alpha}}$, and $B$ may be denoted by
\SourceLeaf{219}
$(ef\ast ke+\alpha)$. If, in place of $R^\mu$, we had considered the term $R^\nu$, we would have obtained a different value of $r$. But it is readily seen that, for any root $\rho$, the root $r$ can have $\frac\lambda{ef}$ different values.

\Needspace{15\baselineskip}\Article{251}
Let us now see how, once one period is labelled, the others are distinguished by their labels. For this purpose, however, we must seek another method of finding the remaining periods: insofar as they, like $A$ itself, are roots of a congruence, no order is discernible among them. Suppose that $A$ is denoted by $(ef\ast0)$. It follows from the preceding discussion that
\[
\begin{aligned}
A^2&\text{ has the form }M+N(ef\ast0)+O(ef\ast1)+P(ef\ast2)+\ldots\\
A^3&\text{ has the form }M'+N'(ef\ast0)+O'(ef\ast1)+\ldots\\
&\text{etc.}\\
A^{ef-1}&\text{ has the form }M^\ast+N^\ast(ef\ast0)+O^\ast(ef\ast1)+\ldots
\end{aligned}
\]
To these is added the congruence
\[
(ef\ast0)+(ef\ast1)+\ldots+(ef\ast ef-1)\equiv-1
\]
We thus have $cf-1$ linear congruences and the same number of unknown quantities, which can therefore be determined by elimination.

\Needspace{6\baselineskip}\emph{Note.} A case may occur in which the unknown quantities are given by expressions of the form $\frac{V}{Wp}$. We shall show below how this difficulty can be remedied. Since this case can occur only very rarely, we do not wish to dwell on it here.

\Article{252}
Let this suffice in general for solving pure congruences. Much more will be said about them at various places below; in particular, many things can be taken over from the solution of pure equations, which we shall not neglect to note in their proper places. We add another example: comparing it with the rules will make everything clearer to those with less experience.

Let $n=31$, $p=311$; that is, we seek the roots of the congruence $x^{31}-1\equiv0\;(\mathrm{mod.}\,311)$. We must first find a primitive root of the congruence $y^{30}-1\equiv\;(\mathrm{mod.}\,31)$, such as $y\equiv3$. Thus put $\rho\equiv3$, and first divide all the primitive roots of the proposed congruence into $5$ periods, namely
\SourceLeaf{220}
\[
\begin{array}{rl}
(5\ast0)\ \dotsfive&(1)+(26)+(25)+(30)+(5)+(6)\\
(5\ast1)\ \dotsfive&(3)+(16)+(13)+(28)+(15)+(18)\\
(5\ast2)\ \dotsfive&(9)+(17)+(8)+(22)+(14)+(23)\\
(5\ast3)\ \dotsfive&(27)+(20)+(24)+(4)+(11)+(7)\\
(5\ast4)\ \dotsfive&(19)+(29)+(10)+(12)+(2)+(21)
\end{array}
\]

The required calculations give the sum of the periods $\equiv-1$, of their squares $\equiv25$, cubes $\equiv26$, fourth powers $\equiv249$, and fifth powers $\equiv564$.

Thus the periods will be the roots of the congruence
\[
x^5+x^4-12x^3-21x^2+x+5\equiv0
\]

Further, one finds
\[
\begin{aligned}
(5\ast0)^2&\equiv6+2(5\ast0)+2(5\ast3)+(5\ast4)\\
(5\ast0)^3&\equiv12+15(5\ast0)+4(5\ast1)+3(5\ast2)+6(5\ast3)+6(5\ast4)\\
(5\ast0)^4&\equiv90+60(5\ast0)+28(5\ast1)+26(5\ast2)+49(5\ast3)+38(5\ast4)
\end{aligned}
\]
and hence, by elimination,
\[
\begin{aligned}
5(5\ast1)&\equiv3(5\ast0)^4-(5\ast0)^3-33(5\ast0)^2-24(5\ast0)+15\\
5(5\ast2)&\equiv-2(5\ast0)^4-(5\ast0)^3+22(5\ast0)^2+31(5\ast0)\\
5(5\ast3)&\equiv(5\ast0)^4-2(5\ast0)^3\\
5(5\ast4)&\equiv-2(5\ast0)^4+4(5\ast0)^3
\end{aligned}
\]

One root of the congruence found above is $\equiv17$. Thus, if $(5\ast0)\equiv17$, then $(5\ast1)\equiv183$, $(5\ast2)\equiv263$, $(5\ast3)\equiv91$, $(5\ast4)\equiv67$.

Now divide each of the periods found into three again, namely
\[
\begin{array}{rl}
(5\ast0)&\text{into }(15\ast0),\ (15\ast5),\ (15\ast10)\\
&\text{or into }(1)+(30),\ (26)+(5),\ (25)+(6)\\
(5\ast1)&\text{into }(15\ast1),\ (15\ast6),\ (15\ast11)\\
&\text{or into }(3)+(28),\ (16)+(15),\ (13)+(18)\\
&\text{etc.}\qquad\text{etc.}
\end{array}
\]

Suppose that the periods into which the following are divided are roots of the respective congruences:
\[
\begin{array}{rll}
(5\ast0)&\text{are roots of}&x^3+Ax^2+Bx+C\equiv0\\
(5\ast1)&&x^3+A'x^2+B'x+C'\equiv0\\
(5\ast2)&&x^3+A''x^2+B''x+C''\equiv0\\
&\text{etc.}
\end{array}
\]

\SourceLeaf{221}
Then
\[
\begin{array}{lll}
A\equiv-(5\ast0),&B\equiv(5\ast0)+(5\ast3),&C\equiv-2-(5\ast4)\\
A'\equiv-(5\ast1),&B'\equiv(5\ast1)+(5\ast4),&C'\equiv-2-(5\ast0)\\
\text{etc.}&\text{etc.}&\text{etc.}
\end{array}
\]

Therefore
\[
\begin{array}{lll}
(15\ast0),\ (15\ast5),\ (15\ast10)&\begin{gathered}\text{will be roots}\\\text{of congruences}\end{gathered}&x^3-17x^2+108x-60\equiv0\\
(15\ast1),\ (15\ast6),\ (15\ast11)&&\\
(15\ast2),\ (15\ast7),\ (15\ast12)&&\\
(15\ast3),\ (15\ast8),\ (15\ast13)&&\\
(15\ast4),\ (15\ast9),\ (15\ast14)&&
\end{array}
\]

Here, however, we have
\[
\begin{gathered}
(15\ast0)^3-3(15\ast0)\equiv(15\ast1)\\
(15\ast1)^3-3(15\ast1)\equiv(15\ast2)\\
\text{etc.}
\end{gathered}
\]

Hence, if one of the roots of the first congruence, $10$, is taken as $(15\ast0)$, we have
\[
\begin{array}{rlrlrl}
(15\ast0)\equiv&10&(15\ast5)\equiv&& (15\ast10)\equiv&\\
(15\ast1)\equiv&37&(15\ast6)\equiv&& (15\ast11)\equiv&\\
(15\ast2)\equiv&-151&(15\ast7)\equiv&& (15\ast12)\equiv&\\
(15\ast3)\equiv&-39&(15\ast8)\equiv&& (15\ast13)\equiv&\\
(15\ast4)\equiv&-112&(15\ast9)\equiv&& (15\ast14)\equiv&
\end{array}
\]

Finally, take the terms making up each of these periods. Then
\[
\begin{array}{lll}
(1),\ (30)&\text{roots of congruences}&x^2-(15\ast0)x+1\equiv0\\
(3),\ (28)&&x^2-(15\ast1)x+1\equiv0\\
&\text{etc.}
\end{array}
\]

The roots of the first congruence are $126$ and $195$, which will therefore be primitive roots of the congruence $x^{31}\equiv1$, and from these the remaining roots can readily be derived.
\BookRule{.28\linewidth}
